\section{Task 1: Extended Experiments and Ablation Studies}
\label{sec:task1_extended}

Following the initial fine-tuning results (Section~\ref{sec:task1}), we conducted several additional experiments to explore alternative approaches suggested in the literature and by collaborators. All experiments are evaluated on the same 166-image test set.

\subsection{Summary of All Approaches}

Table~\ref{tab:all_approaches} summarizes all approaches evaluated for Task 1.

\begin{table}[h]
\centering
\caption{Complete comparison of Task 1 classification approaches}
\label{tab:all_approaches}
\small
\begin{tabular}{lcccc}
\toprule
Approach & Micro P & Micro R & Micro F1 & Note \\
\midrule
Vanilla Florence-2 (zero-shot) & 54.1\% & 18.3\% & 27.4\% & Baseline \\
Fine-tuned v6 (vision+decoder LoRA) & 84.4\% & 62.5\% & 71.9\% & \\
Mixed-data fine-tuning (v9) & 80.2\% & 70.9\% & 75.3\% & \\
Continued fine-tuning (v11) & 83.5\% & 70.5\% & \textbf{76.5\%} & Best \\
v9+v11 ensemble (union) & --- & --- & 76.1\% & Marginal \\
Conservative multi-label (v12b) & --- & --- & 75.9\% & Tied \\
\midrule
\multicolumn{5}{l}{\textit{Alternative approaches (explored but did not improve)}} \\
\midrule
Detection-first (vanilla OD $\rightarrow$ classify) & --- & --- & 52.2\% & \S\ref{ssec:detect_first} \\
Fine-tuned OD labels only & --- & --- & $\sim$40\% & \S\ref{ssec:finetuned_od} \\
Fine-tuned OD $\rightarrow$ crop $\rightarrow$ v11 & 41.7\% & 87.2\% & 56.4\% & \S\ref{ssec:od_crop_classify} \\
OCR input signal (v2) & --- & --- & 48.4\% & \S\ref{ssec:ocr} \\
OCR input signal (v3, conservative) & --- & --- & $\sim$+2\% & \S\ref{ssec:ocr} \\
\bottomrule
\end{tabular}
\end{table}

\subsection{Detection-First Pipeline}
\label{ssec:detect_first}

\textbf{Motivation.} Since under-detection in multi-label images is the main bottleneck, we hypothesized that first detecting individual objects and then classifying each crop might improve recall.

\textbf{Method.} We tested two variants: (1) using vanilla Florence-2's \texttt{<OD>} task for detection followed by v11 classification on each crop, and (2) using a fine-tuned OD model.

\textbf{Vanilla OD $\rightarrow$ v11 classify.} The vanilla Florence-2 object detector produces generic bounding boxes that do not correspond to pantry categories. Combined with v11 classification on crops, this achieved only 52.2\% micro F1, substantially below direct classification.

\subsection{Fine-Tuned Object Detection}
\label{ssec:finetuned_od}

\textbf{Method.} We fine-tuned Florence-2 on the \texttt{<OD>} task using COCO-format bounding box annotations from the training set (10,004 annotations across 21 categories, 15 epochs).

\textbf{Results.} When using the OD model's predicted labels directly as classification output, 7 out of 21 categories achieved 0\% F1. Inspection of the raw output revealed label hallucination (e.g., ``Seafoam -- Fresh'' instead of ``Seafood -- Canned''), indicating that the OD model learned approximate bounding box locations but not stable category labels.

\subsection{Fine-Tuned OD $\rightarrow$ Crop $\rightarrow$ v11 Classify}
\label{ssec:od_crop_classify}

\textbf{Motivation.} Since the fine-tuned OD model produces reasonable bounding boxes even though its labels are unreliable, we combined it with v11: use OD for localization only (ignoring labels), crop each detection, and classify each crop with v11.

\textbf{Method.} The pipeline applies NMS-based deduplication (IoU $> 0.5$), crops each detection with 15\% padding, classifies with v11, and merges with full-image classification as fallback.

\textbf{Results.}

\begin{table}[h]
\centering
\caption{Ablation: OD $\rightarrow$ crop $\rightarrow$ classify pipeline}
\label{tab:od_ablation}
\begin{tabular}{lcccc}
\toprule
Method & Micro P & Micro R & Micro F1 & Macro F1 \\
\midrule
Full image only (v11 baseline) & 83.5\% & 70.5\% & 76.5\% & 75.7\% \\
Crops only & 41.6\% & 83.3\% & 55.5\% & 57.1\% \\
Full image + crops (union) & 41.7\% & 87.2\% & 56.4\% & 57.7\% \\
\bottomrule
\end{tabular}
\end{table}

The crops-based approach dramatically improves recall (70.5\% $\rightarrow$ 87.2\%), confirming that the OD model successfully detects items that v11 misses on the full image. However, precision drops from 83.5\% to 41.7\%, because v11 was trained on full pantry shelf images and over-predicts when given small crops of individual items. The net F1 is substantially lower than direct classification.

\textbf{Key insight:} The detection stage works well for localization, but the classifier is not trained for single-item crops. Retraining the classifier on cropped items could potentially recover precision, but would require generating a new training dataset of individual item crops.

\subsection{OCR-Based Input Signal}
\label{ssec:ocr}

\textbf{Motivation.} Package text (brand names, product descriptions) contains discriminative information that image-only classification may miss. We explored using Florence-2's \texttt{<REGION\_TO\_OCR>} capability as an additional input signal.

\textbf{Method.} Two variants were tested:
\begin{enumerate}[nosep,leftmargin=1.5em]
  \item \textbf{v2 (aggressive):} OCR text matched against a broad keyword list, boosting any matched categories. This was too noisy and \textit{decreased} F1 from 57.6\% to 48.4\%.
  \item \textbf{v3 (conservative):} Added a blacklist for nutrition label text, required multi-word keywords, and required $\geq$2 keyword matches per category. This improved 3 categories (Granola +50\% F1, Savory Snacks +9.4\%, Soup +6.8\%) while minimally impacting others.
\end{enumerate}

\textbf{Conclusion.} OCR as a post-processing signal provides modest gains for select categories but is not a reliable general strategy. The main difficulty is distinguishing informative package text from noise (nutrition labels, legal text, barcodes).

\subsection{Discussion}

Across all alternative approaches, direct generative classification with the fine-tuned Florence-2 model (v11) remained the strongest single method at 76.5\% micro F1. The main findings from the extended experiments are:

\begin{enumerate}[nosep,leftmargin=1.5em]
  \item \textbf{Detection improves recall but destroys precision.} OD-based pipelines can find more items but the classifier over-predicts on small crops, yielding lower overall F1.
  \item \textbf{OCR provides marginal gains.} Text-based signals help for a few categories but introduce too much noise for reliable general use.
  \item \textbf{Generative classification is surprisingly robust.} Despite the multi-label challenge, the generative approach outperforms modular pipelines that decompose the task into detection + classification stages.
  \item \textbf{The remaining bottleneck is data, not architecture.} Class imbalance and limited multi-label training examples constrain performance more than model capacity.
\end{enumerate}

\subsection{Contrastive Learning}
\label{ssec:contrastive}

\textbf{Motivation.} Generative classification requires the model to produce structured JSON output, conflating classification with text generation. We investigated whether a discriminative approach using supervised contrastive learning could decouple feature learning from classification.

\textbf{Method.} We used the frozen Florence-2-large encoder (204M parameters) as a feature extractor, adding two lightweight heads on top of the 1024-dimensional encoder features:
\begin{enumerate}[nosep,leftmargin=1.5em]
  \item \textbf{Projection head} (1024 $\rightarrow$ 256): projects encoder features into a normalized embedding space, trained with supervised contrastive loss (SupCon).
  \item \textbf{Classification head} (1024 $\rightarrow$ 21): multi-label classifier trained with binary cross-entropy loss.
\end{enumerate}
The combined loss is $\mathcal{L} = \alpha \cdot \mathcal{L}_{\text{SupCon}} + (1 - \alpha) \cdot \mathcal{L}_{\text{BCE}}$, with $\alpha = 0.3$ chosen because the SupCon loss magnitude ($\sim$1.4) is much larger than BCE ($\sim$0.05). Only the two heads are trained (\textbf{1.7M trainable parameters}), keeping the 204M encoder parameters frozen. Training used batch size 16, learning rate $5 \times 10^{-4}$ with cosine annealing, temperature $\tau = 0.1$, and a classification threshold of 0.45.

\textbf{Results.}

\begin{table}[h]
\centering
\caption{Contrastive learning variants}
\label{tab:contrastive}
\small
\begin{tabular}{lccccc}
\toprule
Variant & Trainable & Micro P & Micro R & Micro F1 & Macro F1 \\
\midrule
v1 (frozen, $\alpha$=0.5, t=0.07) & 1.7M & 79.2\% & 75.7\% & 77.4\% & 70.9\% \\
v2a (frozen, $\alpha$=0.3, t=0.1) & 1.7M & 82.6\% & 72.1\% & 77.0\% & 69.7\% \\
v2b (partial unfreeze) & $\sim$30M & 72.5\% & 71.3\% & 71.9\% & 64.6\% \\
\midrule
Fine-tuned Florence-2 (v11) & 204M & 82.2\% & 70.1\% & 75.7\% & 74.8\% \\
\bottomrule
\end{tabular}
\end{table}

The contrastive approach produces \textbf{surprisingly competitive results}: the frozen-encoder variant achieves 77.4\% micro F1 with only 0.8\% of the trainable parameters used by the generative approach. Per-class analysis reveals complementary strengths—contrastive excels on categories where generative classification struggles: Ready Meals (84.8\% vs 48.0\%), Savory Snacks (85.7\% vs 64.5\%), and Vegetables Canned (88.2\% vs 69.0\%).

The partially unfrozen variant (v2b) underperforms frozen models, likely due to overfitting on the small training set (166 images). This confirms that the pre-trained Florence-2 encoder features are already well-suited for this task.

\subsection{Ensemble: Generative $\cup$ Contrastive}
\label{ssec:ensemble}

\textbf{Motivation.} Given the complementary strengths of generative (v11) and contrastive classifiers, we explored ensemble strategies to combine both.

\textbf{Method.} Two ensemble strategies were evaluated:
\begin{itemize}[nosep,leftmargin=1.5em]
  \item \textbf{Union ($\cup$):} An item is classified as present if \textit{either} model predicts it. Prioritizes recall.
  \item \textbf{Intersection ($\cap$):} An item is classified as present only if \textit{both} models agree. Prioritizes precision.
\end{itemize}

\textbf{Results.}

\begin{table}[h]
\centering
\caption{Ensemble strategies: generative (v11) $\times$ contrastive}
\label{tab:ensemble}
\begin{tabular}{lcccc}
\toprule
Strategy & Micro P & Micro R & Micro F1 & Macro F1 \\
\midrule
v11 only & 82.2\% & 70.1\% & 75.7\% & 74.8\% \\
Contrastive only & 81.4\% & 74.9\% & 78.0\% & 76.2\% \\
\textbf{Union ($\cup$)} & \textbf{77.2\%} & \textbf{83.7\%} & \textbf{80.3\%} & \textbf{78.4\%} \\
Intersection ($\cap$) & 84.0\% & 66.9\% & 74.5\% & 72.4\% \\
\bottomrule
\end{tabular}
\end{table}

The union ensemble achieves \textbf{80.3\% micro F1}—a +4.6\% absolute improvement over the best single model—by capturing items that one model misses but the other detects. The largest per-class gains are in Ready Meals (+32\%), Savory Snacks (+19.7\%), and Vegetables Canned (+19.3\%). The precision--recall trade-off (77.2\% P / 83.7\% R) favors recall, which is desirable in a food pantry context where missing an item is worse than over-counting.

\subsection{VLM Chain-of-Thought Baseline}
\label{ssec:vlm_cot}

\textbf{Motivation.} Large vision--language models (VLMs) with instruction-following capabilities can perform zero-shot classification through chain-of-thought reasoning, potentially matching or exceeding fine-tuned specialist models without any task-specific training. We evaluate a VLM baseline to contextualize our fine-tuned results.

\textbf{Method.} We used Qwen2.5-VL-7B-Instruct in zero-shot mode with a structured chain-of-thought prompt that instructs the model to: (1) read all text and labels on packaging, (2) identify each food item, and (3) classify into the 21 pantry categories. The prompt includes two few-shot examples and explicit category disambiguation rules. Images were processed at reduced resolution ($256 \times 28^2$ max pixels) using SDPA attention in bf16 precision on a single NVIDIA A100 GPU.

\textbf{Results.}

\begin{table}[h]
\centering
\caption{VLM chain-of-thought zero-shot classification}
\label{tab:vlm_cot}
\begin{tabular}{lcccc}
\toprule
Method & Micro P & Micro R & Micro F1 & Training \\
\midrule
Vanilla Florence-2 (zero-shot) & 54.1\% & 18.3\% & 27.4\% & None \\
Qwen2.5-VL-7B CoT (zero-shot) & 58.3\% & 70.1\% & 63.6\% & None \\
Fine-tuned Florence-2 (v11) & 82.2\% & 70.1\% & 75.7\% & 166 images \\
Ensemble (v11 $\cup$ contrastive) & 77.2\% & 83.7\% & 80.3\% & 166 images \\
\bottomrule
\end{tabular}
\end{table}

The VLM achieves 63.6\% micro F1 in a fully zero-shot setting—more than double the vanilla Florence-2 baseline (27.4\%). Notably, its recall (70.1\%) matches the fine-tuned v11, demonstrating that the VLM's language understanding enables it to identify most items correctly. The gap to the fine-tuned model is primarily in precision (58.3\% vs 82.2\%), as the VLM occasionally hallucinates categories or misinterprets ambiguous packaging.

\textbf{Qualitative analysis.} The VLM's chain-of-thought outputs provide interpretable reasoning. For example, it correctly identifies ``Freshness Honey Iced bread rolls'' from package text but misclassifies them as Bread rather than Desserts—a reasonable interpretation that reveals the inherent ambiguity in some pantry categories. The VLM excels on categories with distinctive packaging (Fresh Fruit: 87.5\% F1, Savory Snacks: 77.8\%) and struggles with visually similar packaged goods (Vegetables Fresh: 28.6\%, Seafood Canned: 40.0\%).

\subsection{Updated Summary}

\begin{table}[h]
\centering
\caption{Complete comparison of all Task 1 approaches}
\label{tab:all_approaches_final}
\small
\begin{tabular}{lcccl}
\toprule
Approach & Micro P & Micro R & Micro F1 & Note \\
\midrule
\multicolumn{5}{l}{\textit{Zero-shot baselines}} \\
\midrule
Vanilla Florence-2 & 54.1\% & 18.3\% & 27.4\% & No training \\
Qwen2.5-VL-7B CoT & 58.3\% & 70.1\% & 63.6\% & No training \\
\midrule
\multicolumn{5}{l}{\textit{Fine-tuned single models}} \\
\midrule
Florence-2 v11 (generative) & 82.2\% & 70.1\% & 75.7\% & Best single (gen.) \\
Contrastive v1 (frozen encoder) & 79.2\% & 75.7\% & 77.4\% & Best single (disc.) \\
Contrastive v2a (tuned hparams) & 82.6\% & 72.1\% & 77.0\% & \\
Contrastive v2b (partial unfreeze) & 72.5\% & 71.3\% & 71.9\% & Overfits \\
\midrule
\multicolumn{5}{l}{\textit{Ensemble}} \\
\midrule
\textbf{v11 $\cup$ contrastive (union)} & \textbf{77.2\%} & \textbf{83.7\%} & \textbf{80.3\%} & \textbf{Best overall} \\
v11 $\cap$ contrastive (intersection) & 84.0\% & 66.9\% & 74.5\% & \\
\midrule
\multicolumn{5}{l}{\textit{Alternative pipelines (did not improve)}} \\
\midrule
OD $\rightarrow$ crop $\rightarrow$ v11 & 41.7\% & 87.2\% & 56.4\% & \\
Detection labels only & --- & --- & $\sim$40\% & \\
OCR input signal (v3) & --- & --- & $\sim$+2\% & Marginal \\
\bottomrule
\end{tabular}
\end{table}
